\documentclass[12pt, a4paper]{article}

\usepackage[utf8]{inputenc}
\usepackage[T1]{fontenc}
\usepackage[french]{babel}
\usepackage{amsmath, amssymb, amsthm}
\usepackage{bm}
\usepackage{graphicx}
\usepackage{hyperref}
\usepackage{geometry}
\usepackage{booktabs}
\usepackage{xcolor}
\usepackage{mathtools}

\geometry{margin=2.5cm}

\hypersetup{
    colorlinks=true,
    linkcolor=blue!70!black,
    citecolor=blue!70!black,
}

\newtheorem{proposition}{Proposition}
\newtheorem{remarque}{Remarque}
\newtheorem{definition}{Définition}

\title{%
    \textbf{Personnalisation de HRTF par apprentissage contrastif cross-modal}\\[0.4em]
    \large Appariement photo d'oreille $\leftrightarrow$ représentation en harmoniques sphériques
}
\author{Mehdi Maqlach}
\date{\today}

% -----------------------------------------------------------------------
\begin{document}
\maketitle

\begin{abstract}
La perception binaurale exacte d'une scène sonore 3D repose sur la
\textit{Head-Related Transfer Function} (HRTF) de l'auditeur, qui code la
signature spectrale imposée par la géométrie individuelle du pavillon, de la
tête et du torse. Les HRTF génériques (moyennées sur une population)
dégradent la localisation, en particulier en élévation, tandis que la mesure
individuelle en chambre anéchoïque reste inaccessible en dehors du
laboratoire. Ce document présente un modèle de \textbf{récupération
cross-modale} : un espace latent partagé est appris entre deux modalités
— une paire de photographies d'oreilles et une HRTF mesurée — de sorte
qu'à l'inférence, la HRTF la plus proche d'un nouvel utilisateur puisse être
retrouvée par similarité cosinus depuis ses seules photos. Nous justifions
le choix d'une représentation en harmoniques sphériques pour la HRTF, une
architecture siamoise pour l'encodeur d'image, une perte contrastive de type
NT-Xent en variante supervisée, et une augmentation de données fondée sur
une symétrie exacte du problème plutôt que sur une interpolation brute.
Le protocole d'évaluation (Recall@$k$, MRR, distorsion log-spectrale) et une
méthode d'interprétabilité par Grad-CAM adaptée à l'architecture sont
également détaillés.
\end{abstract}

\tableofcontents
\vspace{1em}

% =========================================================================
\section{Introduction}

\subsection{Motivation}

Une HRTF $H(\varphi,\theta,f)$ décrit le filtrage acoustique subi par une
onde sonore provenant de la direction $(\varphi,\theta)$ avant d'atteindre
le tympan. Ce filtrage dépend fortement de la morphologie individuelle du
pavillon — ses replis (hélix, anthélix, conque) créent des résonances et
des annulations spécifiques à chaque sujet, en particulier au-dessus de
6\,kHz, qui sont les indices dominants de la perception de l'élévation et de
la résolution avant/arrière.

Deux stratégies existent pour spatialiser un signal en binaural~:

\begin{itemize}
    \item convoluer avec une HRTF \textbf{générique} (moyenne de population)
    — robuste mais acoustiquement fausse pour tout individu s'écartant de
    la moyenne, source de confusions avant/arrière et d'une élévation
    perçue aplatie~;
    \item convoluer avec la HRTF \textbf{individuelle}, mesurée en chambre
    anéchoïque — fidèle mais nécessitant un protocole de mesure lourd
    (réseau de haut-parleurs, immobilisation du sujet), inapplicable à
    grande échelle.
\end{itemize}

\subsection{Approche retenue}

Nous formulons la personnalisation non pas comme un problème de
\emph{régression} (prédire directement les échantillons temporels d'une
HRTF depuis une image, tâche mal posée avec un nombre de sujets mesurés
nécessairement restreint), mais comme un problème de \textbf{retrieval
cross-modal}~: apprendre une fonction d'encodage pour chaque modalité telle
que les deux représentations d'un même sujet soient proches dans un espace
latent commun, puis retrouver par plus-proche-voisin la HRTF mesurée la
plus vraisemblable pour un nouvel utilisateur.

Cette formulation, popularisée en vision-langage par CLIP~\cite{Radford21},
présente un avantage décisif pour notre contexte~: la sortie du modèle est
toujours une HRTF \emph{physiquement valide} (une mesure réelle de la base
de référence), ce qui élimine le risque de synthétiser un filtre non
causal ou acoustiquement aberrant — risque inhérent à toute régression
directe entraînée sur peu d'exemples.

\subsection{Contributions}

\begin{enumerate}
    \item Une représentation de la HRTF en \textbf{harmoniques sphériques}
    (SH), de taille fixe et indépendante de la grille de mesure du sujet
    (détail complet dans le document compagnon~\cite{SHCompanion}),
    utilisée ici comme entrée de l'encodeur HRTF.
    \item Une architecture à deux encodeurs (\texttt{EarEncoder},
    \texttt{HRTFEncoder}) projetant chaque modalité dans un espace latent
    $L_2$-normalisé partagé.
    \item Une perte contrastive \textbf{NT-Xent supervisée}, choisie pour
    traiter correctement les échantillons augmentés d'un même sujet comme
    des positifs multiples plutôt que comme des faux négatifs.
    \item Une augmentation de données par \textbf{réflexion azimutale
    exacte}, dérivée analytiquement de la parité des harmoniques
    sphériques réelles — et non d'une interpolation ou d'un bruit ajouté.
    \item Un protocole d'entraînement en trois phases (\textit{progressive
    unfreezing}) et un protocole d'évaluation anti-fuite adapté à un
    nombre de sujets restreint.
\end{enumerate}

% =========================================================================
\section{Travaux connexes}

\textbf{Personnalisation de HRTF.} Les approches historiques reposent sur
des mesures anthropométriques (largeur de tête, taille du pavillon) reliées
à la HRTF par régression ou par sélection dans une base
(\textit{best-matched HRTF}). Notre approche s'en distingue en utilisant
directement l'image plutôt que des mesures manuelles, et en formulant le
problème comme un retrieval plutôt qu'une régression paramétrique.

\textbf{Apprentissage contrastif.} SimCLR~\cite{Chen19} et la perte
InfoNCE~\cite{Oord18} (elle-même héritière de la N-pair loss~\cite{Sohn16})
établissent le cadre général~: rapprocher les représentations de paires
positives et éloigner celles de paires négatives via une cross-entropie
softmax température-contrôlée. CLIP~\cite{Radford21} étend ce cadre au
cas \textbf{cross-modal} (image $\leftrightarrow$ texte) — c'est le cadre
que nous transposons ici à (image d'oreille $\leftrightarrow$ HRTF). La
\textit{Supervised Contrastive Loss}~\cite{Khosla20} généralise encore ce
cadre en autorisant plusieurs positifs par ancre, ce qui est notre cas en
présence d'augmentation par sujet.

\textbf{Interprétabilité.} Grad-CAM~\cite{Selvaraju17} est repris ici dans
une variante adaptée à une tête d'embedding contrastive $L_2$-normalisée
(section~\ref{sec:gradcam}), un cas non standard où le score de classe
usuel n'existe pas.

% =========================================================================
\section{Données et représentation}

\subsection{Deux bases de données, deux rôles}

Le projet s'appuie sur deux corpus HRTF distincts, à des fins différentes~:

\begin{center}
\begin{tabular}{lll}
\toprule
Base & Rôle & Contenu \\
\midrule
IRCAM LISTEN & moteur DSP (interpolation, synthèse) & 6 sujets, grille de 187 positions \\
Base cross-modale\footnotemark & entraînement du modèle contrastif & sujets avec HRIR + photos d'oreilles G/D \\
\bottomrule
\end{tabular}
\end{center}
\footnotetext{Structure compatible HUTUBS~\cite{Brinkmann19} / SONICOM~:
un dossier par sujet contenant un fichier \texttt{*\_SOFA.sofa} et des
scans/photos d'oreille gauche et droite.}

Le modèle décrit dans ce document est entraîné exclusivement sur la seconde
base. Une exécution observée dans le pipeline d'entraînement donne, après
expansion par augmentation (section~\ref{sec:augmentation}), un total de
$80$ sujets répartis en $57$ (train) / $11$ (val) / $12$ (test), avec un
ordre SH $L=10$ (soit $121$ coefficients par oreille).

\subsection{Décomposition en harmoniques sphériques}
\label{sec:sh}

Les positions de mesure varient d'un sujet à l'autre et d'une base à
l'autre (de $187$ à plusieurs milliers de points), ce qui interdit de
représenter la HRTF comme un tenseur de forme fixe sans prétraitement. La
solution retenue projette la HRTF sur la base orthonormée des harmoniques
sphériques réelles $Y_l^m$ :

\begin{equation}
    H(\varphi,\theta,f) \;\approx\;
    \sum_{l=0}^{L}\sum_{m=-l}^{l} a_l^m(f)\, Y_l^m(\varphi,\theta),
    \qquad
    Y_l^m(\varphi,\theta) =
    \begin{cases}
        \sqrt{2}\,K_l^m \cos(m\varphi)\,P_l^{m}(\cos\theta) & m>0\\
        K_l^0\,P_l^0(\cos\theta) & m=0\\
        \sqrt{2}\,K_l^{|m|} \sin(|m|\varphi)\,P_l^{|m|}(\cos\theta) & m<0
    \end{cases}
    \label{eq:sh-basis}
\end{equation}

avec $K_l^m = \sqrt{\tfrac{2l+1}{4\pi}\tfrac{(l-|m|)!}{(l+|m|)!}}$. Pour un
ordre de troncature $L$, le nombre de coefficients est fixe et vaut
$n_{sh}=(L+1)^2$ — $121$ pour $L=10$ — indépendamment du nombre de
positions mesurées.

Les coefficients sont obtenus par moindres carrés~: en notant
$\mathbf{Y}\in\mathbb{R}^{N\times n_{sh}}$ la matrice de design
($\mathbf{Y}_{i,lm}=Y_l^m(\varphi_i,\theta_i)$) et
$\mathbf{h}(f)\in\mathbb{R}^N$ le vecteur de log-magnitudes mesurées à la
fréquence $f$,

\begin{equation}
    \hat{\mathbf{a}}(f) = \mathbf{Y}^{+}\mathbf{h}(f)
    = \bigl(\mathbf{Y}^\top\mathbf{Y}\bigr)^{-1}\mathbf{Y}^\top \mathbf{h}(f).
    \label{eq:sh-lstsq}
\end{equation}

Appliqué à chaque bin fréquentiel et aux deux oreilles, ceci produit le
tenseur $\mathbf{A}\in\mathbb{R}^{n_{sh}\times N_f \times 2}$ qui constitue
l'entrée de l'encodeur HRTF ($N_f=129$ pour un FIR de $256$ échantillons à
$44{,}1$\,kHz). Cette représentation est justifiée en détail — orthogonalité,
complétude, hiérarchie spatiale interprétable, lien avec les modes propres
de l'équation de Laplace sur $S^2$ — dans le document compagnon
\textit{Représentation des HRTFs par Harmoniques Sphériques}~\cite{SHCompanion};
nous n'en reprenons ici que les éléments strictement nécessaires à
l'architecture du modèle.

\subsection{Prétraitement des images}

Chaque photo d'oreille est redimensionnée en $224\times224\times3$
(interpolation Lanczos) et normalisée dans $[0,1]$. Aucune normalisation
supplémentaire n'est appliquée en amont du réseau~: le \textit{rescaling}
vers l'échelle attendue par EfficientNet ($[0,255]$) est intégré comme
première couche du modèle (section~\ref{sec:ear-encoder}), afin que la
même fonction de prétraitement s'applique identiquement à l'entraînement
et à l'inférence.

Les coefficients SH sont normalisés en $z$-score,
$\tilde{a} = (a-\mu_{\mathrm{train}})/\sigma_{\mathrm{train}}$, avec
$\mu_{\mathrm{train}}, \sigma_{\mathrm{train}}$ estimés uniquement sur le
split d'entraînement — condition nécessaire pour ne pas laisser fuiter
d'information des sujets de validation/test dans la normalisation.

% =========================================================================
\section{Formulation du problème}

\begin{definition}[Espace latent partagé]
Soit $I_L, I_R$ les photos d'oreille gauche/droite d'un sujet et
$\mathbf{A}$ ses coefficients SH. On cherche deux fonctions
$$
f_{\mathrm{ear}} : (I_L,I_R) \longmapsto z_{\mathrm{ear}} \in \mathbb{S}^{d-1},
\qquad
f_{\mathrm{hrtf}} : \mathbf{A} \longmapsto z_{\mathrm{hrtf}} \in \mathbb{S}^{d-1},
$$
paramétrées par des réseaux de neurones et projetant sur la sphère unité de
$\mathbb{R}^d$ ($d=128$), telles que pour un même sujet $i$, $z_{\mathrm{ear},i}$
et $z_{\mathrm{hrtf},i}$ soient les plus proches possibles au sens du
cosinus, et éloignés de ceux de tout autre sujet.
\end{definition}

La contrainte $z\in\mathbb{S}^{d-1}$ ($L_2$-normalisation en sortie des deux
encodeurs) n'est pas cosmétique~: elle rend le produit scalaire équivalent
à la similarité cosinus,
$\langle z_{\mathrm{ear}}, z_{\mathrm{hrtf}}\rangle = \cos(z_{\mathrm{ear}}, z_{\mathrm{hrtf}})$,
condition requise par la perte NT-Xent (section~\ref{sec:loss}) et par la
recherche du plus proche voisin en inférence (section~\ref{sec:inference}).

% =========================================================================
\section{Architecture du modèle}

\subsection{EarEncoder — réseau siamois}
\label{sec:ear-encoder}

\begin{equation}
\begin{aligned}
    I_L,I_R &\in [0,1]^{224\times224\times3} \\
    F_L = \mathrm{EffNetB0}(255\cdot I_L), \quad F_R &= \mathrm{EffNetB0}(255\cdot I_R)
    \qquad \text{(poids \textbf{partagés})} \\
    x &= \mathrm{Concat}\bigl[\mathrm{GAP}(F_L),\,\mathrm{GAP}(F_R)\bigr] \in \mathbb{R}^{2560}\\
    x &\to \mathrm{Dense}(512,\mathrm{relu}) \to \mathrm{Dropout}(0.3)
        \to \mathrm{Dense}(256,\mathrm{relu}) \to \mathrm{Dense}(128) \\
    z_{\mathrm{ear}} &= x / \lVert x \rVert_2 \in \mathbb{S}^{127}
\end{aligned}
\end{equation}

Le partage des poids entre les deux branches gauche/droite (\textit{Siamese
network}) n'est pas un simple gain de paramètres~: il encode a priori que
la fonction \textit{morphologie du pavillon} $\to$ \textit{signature
spectrale} est la même quelle que soit l'oreille, et réduit d'un facteur
deux la variance d'estimation pour un nombre de sujets mesurés restreint.

Le \textit{backbone} EfficientNet-B0~\cite{Tan19} est pré-entraîné sur
ImageNet. Le \textit{fine-tuning} se fait en trois phases de dégel
progressif (section~\ref{sec:training}), une stratégie standard en
apprentissage par transfert avec peu de données~: entraîner d'abord la tête
seule évite que des gradients de large amplitude, issus d'une tête
initialisée aléatoirement, ne détruisent les features génériques du
backbone pré-entraîné.

\subsection{HRTFEncoder}

\begin{equation}
\begin{aligned}
    \mathbf{A} &\in \mathbb{R}^{n_{sh}\times N_f \times 2}
    \;\xrightarrow{\text{reshape}}\;
    \mathbb{R}^{n_{sh}\times (2N_f)} \\
    &\xrightarrow{\mathrm{Dense}(32,\mathrm{relu})\ \text{partagé sur } n_{sh}}\;
    \mathbb{R}^{n_{sh}\times 32}
    \;\xrightarrow{\text{flatten}}\;
    \mathbb{R}^{32\, n_{sh}} \\
    &\to \mathrm{Dense}(256,\mathrm{relu}) \to \mathrm{Dropout}(0.3)
        \to \mathrm{Dense}(128)
    \;\to\; z_{\mathrm{hrtf}} = x/\lVert x\rVert_2 \in \mathbb{S}^{127}
\end{aligned}
\end{equation}

La couche \texttt{Dense(32)} est appliquée identiquement à chacune des
$n_{sh}$ lignes de coefficients avant l'aplatissement — elle joue le rôle
d'un encodeur \emph{par harmonique}, cohérent avec le fait que chaque ligne
$a_l^m(\cdot)$ est une composante spatiale décorrélée des autres (propriété
d'orthogonalité, éq.~\eqref{eq:sh-basis}).

% =========================================================================
\section{Fonction de perte}
\label{sec:loss}

\subsection{NT-Xent standard}

Pour un batch de $N$ paires $(z_{\mathrm{ear},i}, z_{\mathrm{hrtf},i})_{i=1}^N$,
on définit la similarité température-contrôlée
$s_{ij} = \langle z_{\mathrm{ear},i}, z_{\mathrm{hrtf},j}\rangle / \tau$.
Le problème de retrouver la HRTF du sujet $i$ parmi les $N$ candidats du
batch est traité comme une classification à $N$ classes~:

\begin{equation}
    \mathcal{L}_{\mathrm{ear}\to\mathrm{hrtf}}
    = -\frac{1}{N}\sum_{i=1}^{N} \log
    \frac{\exp(s_{ii})}{\sum_{k=1}^N \exp(s_{ik})},
    \qquad
    \mathcal{L}_{\mathrm{hrtf}\to\mathrm{ear}}
    = -\frac{1}{N}\sum_{i=1}^{N} \log
    \frac{\exp(s_{ii})}{\sum_{k=1}^N \exp(s_{ki})}
\end{equation}

\begin{equation}
    \boxed{\;\mathcal{L}_{\mathrm{NT\text{-}Xent}}
    = \tfrac{1}{2}\bigl(\mathcal{L}_{\mathrm{ear}\to\mathrm{hrtf}}
    + \mathcal{L}_{\mathrm{hrtf}\to\mathrm{ear}}\bigr)\;}
\end{equation}

La symétrisation force l'espace latent à être cohérent dans les deux sens
de récupération (oreille $\to$ HRTF \emph{et} HRTF $\to$ oreille), plutôt
que de sur-optimiser un seul sens au détriment de l'autre.

\subsection{Supervised NT-Xent — nécessaire en présence d'augmentation}

L'augmentation par réflexion azimutale (section~\ref{sec:augmentation})
crée, pour un sujet de base $b$, plusieurs échantillons
$(b, b_{\mathrm{mirror}}, \dots)$ dans le même batch. Sous NT-Xent
standard, deux échantillons du même sujet de base occupant des lignes
différentes de la matrice $s$ seraient traités comme des \textbf{négatifs}
l'un de l'autre — un faux négatif structurel qui pousse le modèle à
\emph{séparer} des représentations qui devraient au contraire être
proches.

Soit $P(i) = \{\,j : \mathrm{label}(j) = \mathrm{label}(i)\,\}$ l'ensemble
des indices du batch partageant le même sujet de base que $i$ (labels
attribués via l'identifiant avant suffixe, \emph{ex.} \texttt{H10\_mirror}
$\to$ \texttt{H10}). La perte supervisée~\cite{Khosla20}, adaptée au cadre
cross-modal, moyenne la log-vraisemblance sur tous les positifs de
l'ancre~:

\begin{equation}
    \mathcal{L}_i^{\mathrm{ear}\to\mathrm{hrtf}}
    = -\frac{1}{|P(i)|}\sum_{p\in P(i)}
    \log \frac{\exp(s_{ip})}{\sum_{k=1}^N \exp(s_{ik})}
    \label{eq:sup-ntxent}
\end{equation}

et symétriquement pour $\mathcal{L}_i^{\mathrm{hrtf}\to\mathrm{ear}}$; la
perte totale est la moyenne batch des deux termes sur $i=1,\dots,N$. C'est
ce mode (\texttt{supervised=True}) qui est utilisé par défaut dès lors que
le dataset est augmenté.

\subsection{Température $\tau$}

$\tau=0{,}07$ (valeur reprise de SimCLR~\cite{Chen19}) contrôle la netteté
de la distribution softmax sur les similarités~: un $\tau$ faible pénalise
fortement les négatifs proches du positif (utile quand $N$ est petit, comme
ici, pour extraire un signal d'apprentissage suffisant d'un nombre limité
de négatifs par batch), tandis qu'un $\tau$ élevé lisse la distribution et
ralentit la convergence.

% =========================================================================
\section{Augmentation par symétrie azimutale}
\label{sec:augmentation}

\subsection{Motivation}

Avec seulement quelques dizaines de sujets mesurés, l'augmentation de
données est indispensable. Une interpolation ou un \textit{mixup} naïf
entre deux HRTF de sujets différents introduit des artefacts de phase
(deux HRIR aux délais d'onset différents s'annulent partiellement en
moyenne complexe directe) sans justification acoustique. On lui préfère
ici une transformation qui est une \textbf{symétrie exacte} du problème~:
la réflexion gauche/droite d'une tête (approximativement) bilatéralement
symétrique.

\subsection{Dérivation}

\begin{proposition}[Symétrie azimutale des coefficients SH réels]
Soit un sujet de HRTF gauche $H_L(\varphi,\theta,f)$ et droite
$H_R(\varphi,\theta,f)$. Sous l'hypothèse de symétrie bilatérale de la
tête, le sujet miroir obtenu en réfléchissant l'azimut ($\varphi \to
-\varphi$) satisfait
$H_L^{\mathrm{mirror}}(\varphi,\theta,f) = H_R(-\varphi,\theta,f)$ et
réciproquement. En notant $a_l^m$ les coefficients SH du sujet original,
les coefficients du sujet miroir vérifient
\begin{equation}
    a_l^{m,\,\mathrm{mirror}} = \sigma(m)\cdot a_l^{m},
    \qquad
    \sigma(m) = \begin{cases} -1 & m < 0 \\ +1 & m \geq 0 \end{cases}
    \label{eq:mirror-sign}
\end{equation}
avec échange simultané des canaux gauche/droite.
\end{proposition}

\begin{proof}
Dans l'éq.~\eqref{eq:sh-basis}, la dépendance en $\varphi$ n'apparaît que
via $\cos(m\varphi)$ (termes $m>0$) et $\sin(|m|\varphi)$ (termes $m<0$);
le facteur de Legendre $P_l^{|m|}(\cos\theta)$ est invariant par
$\varphi\to-\varphi$. Or $\cos(-m\varphi)=\cos(m\varphi)$ (fonction paire)
tandis que $\sin(-|m|\varphi) = -\sin(|m|\varphi)$ (fonction impaire).
Comme la projection SH (éq.~\eqref{eq:sh-lstsq}) est linéaire en les
valeurs de $H$, la réflexion des positions se traduit exactement par le
changement de signe~\eqref{eq:mirror-sign} appliqué aux coefficients, sans
qu'aucune ré-estimation par moindres carrés ne soit nécessaire.
\end{proof}

\begin{remarque}
Cette transformation est \textbf{exacte et sans perte}, contrairement à
une augmentation stochastique (bruit gaussien, \textit{mixup} linéaire)~:
elle ne fait qu'exploiter une symétrie déjà présente dans la base $Y_l^m$.
C'est pour cette même raison que l'interpolation phase/onset — et non une
interpolation complexe directe — est utilisée ailleurs dans le projet
(module d'interpolation spatiale de la HRTF) : dans les deux cas, on évite
d'introduire un artefact de phase en respectant la structure du problème
plutôt qu'en interpolant des représentations où cette structure n'est pas
linéaire.
\end{remarque}

En parallèle, les images sont transformées par échange gauche/droite suivi
d'un miroir horizontal~: $I_L^{\mathrm{mirror}} = \mathrm{flip}(I_R)$,
$I_R^{\mathrm{mirror}} = \mathrm{flip}(I_L)$ — cohérent géométriquement
avec la réflexion appliquée aux coefficients SH.

\subsection{Application offline et online, anti-fuite}

L'expansion $N \to 2N$ sujets (\texttt{expand()}) est appliquée
\textbf{avant} le split train/val/test, afin que chaque sujet miroir soit
disponible pour la même expansion en aval. Une seconde couche
d'augmentation, appliquée \textit{online} par \texttt{tf.data.map()}
\textbf{uniquement sur le train set}, permet de tirer une réflexion
aléatoire (probabilité $p=0{,}5$) à chaque epoch — le val/test set restent
non augmentés pour donner une mesure honnête de la performance.

Le fait que l'expansion précède le split introduit malgré tout une fuite
partielle~: un sujet peut se retrouver en version originale dans le train
set et en version miroir dans le val set. Le protocole d'évaluation
(section~\ref{sec:evaluation}) neutralise cette fuite en comparant les
identifiants de \emph{sujet de base} plutôt que les identifiants exacts.

% =========================================================================
\section{Entraînement}
\label{sec:training}

\subsection{Dégel progressif en trois phases}

\begin{center}
\begin{tabular}{cccc}
\toprule
Phase & Backbone EfficientNet-B0 & Paramètres entraînables & Taux d'apprentissage \\
\midrule
1 & entièrement gelé & $\sim 1{,}4$\,M (tête seule) & $10^{-3}$ \\
2 & 20 dernières couches dégelées & $\sim 5$\,M & $10^{-4}$ \\
3 & entièrement dégelé & $\sim 6{,}5$\,M & $10^{-5}$ \\
\bottomrule
\end{tabular}
\end{center}

Cette décroissance du taux d'apprentissage de deux ordres de grandeur au
fil des phases est le pendant nécessaire de l'ouverture progressive du
nombre de paramètres entraînables~: un $\mathrm{lr}=10^{-3}$ appliqué au
backbone entier (phase 3) détruirait les poids pré-entraînés avant que la
tête n'ait eu le temps de fournir un signal de gradient informatif.

\subsection{Callbacks et infrastructure}

Chaque phase est équipée de \texttt{EarlyStopping} (patience $10$,
restauration des meilleurs poids), \texttt{ReduceLROnPlateau} (facteur
$0{,}5$, patience $5$) et \texttt{ModelCheckpoint}. Un
\texttt{RecallCallback} recalcule, à chaque epoch, une gallery
d'embeddings HRTF depuis le train set et mesure le Recall@1 des photos du
val set (section~\ref{sec:evaluation}) — une métrique de bout en bout, plus
directement interprétable que la seule valeur de la perte. L'ensemble est
tracé et versionné via TensorBoard et MLflow (hash du dataset, hyperparamètres,
courbes, artefacts de poids par phase), pour garantir la reproductibilité
des expériences.

\subsection{Extension~: apprentissage fédéré}

Les photos d'oreilles sont une donnée biométrique sensible. Une extension
du \texttt{Trainer} vers un cadre fédéré (\textbf{FedAvg}~\cite{McMahan17},
via le framework Flower) permet à plusieurs clients d'entraîner localement
le \texttt{ContrastiveModel} sur leurs propres sujets sans jamais
centraliser les images~: seuls les poids du modèle (ou leurs mises à jour,
potentiellement chiffrées de façon homomorphe pour l'agrégation) transitent
vers le serveur. Cette direction, explorée dans le projet à titre de preuve
de concept, est discutée plus en détail en perspective (section~\ref{sec:discussion}).

% =========================================================================
\section{Inférence}
\label{sec:inference}

\subsection{Gallery}

La \texttt{Gallery} est une base de référence construite \textbf{une seule
fois} depuis le train set~: pour chaque sujet connu, elle stocke
l'embedding $z_{\mathrm{hrtf}}$ ($L_2$-normalisé) et les coefficients SH
dénormalisés associés. La recherche du plus proche voisin est un simple
produit scalaire,

\begin{equation}
    \mathrm{score}(i) = \langle z_{\mathrm{hrtf},i},\, z_{\mathrm{ear}} \rangle,
    \qquad
    \hat{\imath} = \arg\max_i\, \mathrm{score}(i),
\end{equation}

les deux membres étant sur $\mathbb{S}^{127}$, ce score est directement un
cosinus $\in[-1,1]$.

\subsection{HRTFPredictor}

Depuis deux photos d'un nouvel utilisateur, le pipeline complet retourne le
sujet le plus proche, son score de similarité, et ses coefficients SH — que
l'on peut ensuite \textbf{décoder à n'importe quelle direction} via
$\hat{H}(\varphi,\theta,f)=\sum_{l,m} a_l^m(f)\,Y_l^m(\varphi,\theta)$
(éq.~\eqref{eq:sh-basis}), pour alimenter le moteur de convolution binaurale
du projet (\texttt{HRTFInterpolator}, \texttt{DynamicConvolver}).

% =========================================================================
\section{Évaluation}
\label{sec:evaluation}

\subsection{Métriques de retrieval}

Pour un ensemble de $N$ sujets test avec matrice de similarité
$s_{ij}=\langle z_{\mathrm{ear},i}, z_{\mathrm{hrtf},j}\rangle$~:

\begin{equation}
    \mathrm{Recall@}k = \frac{1}{N}\sum_{i=1}^N
    \mathbb{1}\bigl[i \in \mathrm{top}_k(s_{i,\cdot})\bigr],
    \qquad
    \mathrm{MRR} = \frac{1}{N}\sum_{i=1}^N \frac{1}{\mathrm{rang}(i)}
\end{equation}

où $\mathrm{rang}(i)$ est la position du bon sujet dans le classement des
similarités décroissantes. Ces deux métriques mesurent directement la tâche
opérationnelle (retrouver la bonne HRTF), indépendamment de l'échelle de
la perte contrastive.

\subsection{Distorsion log-spectrale (LSD)}

Pour comparer la HRTF récupérée à la HRTF réelle du sujet, on calcule la
distorsion log-spectrale sur les coefficients SH,

\begin{equation}
    \mathrm{LSD} = \sqrt{\,\overline{\bigl(20\log_{10}|A_{\mathrm{pred}}|
    - 20\log_{10}|A_{\mathrm{true}}|\bigr)^2}\,} \quad \text{[dB]},
\end{equation}

la moyenne étant prise sur tous les coefficients SH, bins fréquentiels et
canaux. Une valeur inférieure à $\sim5$\,dB est généralement considérée
comme acoustiquement acceptable dans la littérature HRTF.

\subsection{Séparabilité de l'espace latent}

Le rapport entre similarité moyenne des paires correctes ($S^+$) et des
paires incorrectes ($S^-$) donne une mesure globale de la structure de
l'espace appris — une valeur cible $S^+/S^- > 1{,}5$ indique une séparation
nette entre sujets plutôt qu'un espace latent effondré.

\subsection{Protocole anti-fuite}

Comme évoqué en section~\ref{sec:augmentation}, l'expansion par réflexion
précède le split. Le critère de succès des métriques ci-dessus est donc
appliqué sur l'\textbf{identifiant de sujet de base}
(\texttt{H10\_mirror}\,$\to$\,\texttt{H10}) et non sur l'identifiant exact,
pour éviter de compter comme une fausse victoire une reconnaissance qui ne
ferait qu'exploiter la fuite miroir train/val. Il s'agit néanmoins d'une
évaluation \emph{optimiste}~: un split réalisé avant augmentation serait
plus rigoureux, et constitue une amélioration identifiée pour une itération
future du protocole.

% =========================================================================
\section{Interprétabilité}
\label{sec:gradcam}

L'architecture pose un problème non standard pour Grad-CAM~\cite{Selvaraju17}~:
la sortie est un embedding $L_2$-normalisé, pas un score de classe. Le
score usuel $\lVert z \rVert_2^2 / d = 1/d$ est \textbf{constant} sur la
sphère unité — son gradient par rapport aux cartes d'activation est
identiquement nul, produisant des heatmaps vides. On lui substitue un score
$\sum_j |z_j|$ (norme $\ell_1$, non constante sur $\mathbb{S}^{d-1}$), ou,
pour expliquer \emph{pourquoi} une paire (oreille, HRTF) a été mise en
correspondance, le score de similarité
$\mathrm{score}(z) = \langle z,\, z_{\mathrm{hrtf}}^{\mathrm{ref}}\rangle$
avec un embedding HRTF cible fixé.

Une seconde difficulté est spécifique au backbone gelé (phases 1--2)~: un
appel monolithique du modèle Keras à travers une \texttt{GradientTape} ne
retourne aucun gradient vers les cartes de features, la tape ne surveillant
pas les variables non-entraînables et l'appel du modèle constituant une
frontière opaque au graphe de gradient. La solution consiste à séparer
explicitement le calcul en deux étapes~: le backbone est exécuté \emph{hors}
tape (sa trainabilité devient alors indifférente, on ne demande jamais de
gradient par rapport à ses poids), puis la tête est appelée \emph{couche
par couche} à l'intérieur de la tape avec \texttt{tape.watch()} explicite
sur les cartes de features en sortie du backbone.

% =========================================================================
\section{Discussion et limites}
\label{sec:discussion}

\begin{itemize}
    \item \textbf{Taille du dataset.} Quelques dizaines de sujets mesurés
    (avant augmentation) restent peu nombreux pour un apprentissage
    contrastif, dont la qualité dépend directement du nombre de négatifs
    disponibles par batch. L'augmentation par symétrie exacte
    (section~\ref{sec:augmentation}) atténue mais ne supprime pas cette
    contrainte.
    \item \textbf{Fuite partielle train/val.} Le protocole d'anti-fuite par
    identifiant de base (section~\ref{sec:evaluation}) reste une mesure
    optimiste; un split \emph{avant} expansion offrirait une estimation de
    généralisation plus fiable, au prix d'un dataset d'entraînement plus
    petit.
    \item \textbf{Biais de la base de mesure.} Toute base HRTF mesurée
    sur-représente certaines morphologies selon la population recrutée; la
    Gallery ne peut retrouver que des sujets présents dans sa base de
    référence — un nouvel utilisateur à la morphologie atypique reçoit la
    HRTF la plus proche disponible, pas une HRTF réellement personnalisée.
    \item \textbf{Qualité de l'image.} La performance du modèle dépend de
    conditions de prise de vue (éclairage, angle, cadrage) non contrôlées
    en usage réel, contrairement aux scans standardisés du dataset
    d'entraînement.
\end{itemize}

% =========================================================================
\section{Conclusion et perspectives}

Ce document formalise un modèle de personnalisation de HRTF par
récupération cross-modale, où la HRTF est représentée de façon compacte et
physiquement motivée par ses coefficients en harmoniques sphériques, et où
l'alignement avec une paire de photos d'oreilles est appris par perte
contrastive supervisée dans un espace latent $L_2$-normalisé partagé. Le
choix de formuler le problème comme un retrieval plutôt qu'une régression,
l'augmentation par symétrie exacte plutôt que par interpolation, et le
dégel progressif du backbone, répondent tous à la même contrainte
structurante du projet~: un nombre de sujets mesurés intrinsèquement
limité.

Perspectives directes~: (i) évaluation avec split réalisé avant
augmentation, pour une mesure de généralisation non optimiste; (ii)
extension de l'apprentissage fédéré à un cadre multi-institutionnel réel,
avec agrégation chiffrée; (iii) fusion, à l'inférence, du sujet retrouvé
par Gallery avec le moteur d'interpolation spatiale du projet, pour
produire une HRTF personnalisée continue en direction plutôt que limitée
aux positions mesurées du sujet retrouvé.

% -----------------------------------------------------------------------
\begin{thebibliography}{99}

\bibitem{SHCompanion}
M.~Maqlach,
\textit{Représentation des HRTFs par Harmoniques Sphériques — Formalisme
mathématique et justification pour l'encodage cross-modal},
document technique, \texttt{docs/hrtf\_spherical\_harmonics.tex}, 2026.

\bibitem{Evans98}
M.~A. Evans, J.~A. S. Angus, and A.~I. Tew,
\textit{Analyzing head-related transfer function measurements using
surface spherical harmonics},
J.\ Acoust.\ Soc.\ Am., 104(4):2400--2411, 1998.

\bibitem{Zotkin09}
D.~N. Zotkin, R.~Duraiswami, and N.~A. Gumerov,
\textit{Regularized HRTF fitting using spherical harmonics},
in Proc.\ IEEE WASPAA, 2009.

\bibitem{Chen19}
T.~Chen, S.~Kornblith, M.~Norouzi, and G.~Hinton,
\textit{A simple framework for contrastive learning of visual
representations (SimCLR)},
ICML, 2020.

\bibitem{Oord18}
A.~van den Oord, Y.~Li, and O.~Vinyals,
\textit{Representation learning with contrastive predictive coding},
arXiv:1807.03748, 2018.

\bibitem{Sohn16}
K.~Sohn,
\textit{Improved deep metric learning with multi-class N-pair loss
objective},
NeurIPS, 2016.

\bibitem{Khosla20}
P.~Khosla, P.~Teterwak, C.~Wang, A.~Sarna, Y.~Tian, P.~Isola, A.~Maschinot,
C.~Liu, and D.~Krishnan,
\textit{Supervised contrastive learning},
NeurIPS, 2020.

\bibitem{Radford21}
A.~Radford, J.~W. Kim, C.~Hallacy, A.~Ramesh, G.~Goh, S.~Agarwal,
G.~Sastry, A.~Askell, P.~Mishkin, J.~Clark, G.~Krueger, and I.~Sutskever,
\textit{Learning transferable visual models from natural language
supervision (CLIP)},
ICML, 2021.

\bibitem{Tan19}
M.~Tan and Q.~Le,
\textit{EfficientNet: Rethinking model scaling for convolutional neural
networks},
ICML, 2019.

\bibitem{Selvaraju17}
R.~R. Selvaraju, M.~Cogswell, A.~Das, R.~Vedantam, D.~Parikh, and D.~Batra,
\textit{Grad-CAM: Visual explanations from deep networks via
gradient-based localization},
ICCV, 2017.

\bibitem{McMahan17}
H.~B. McMahan, E.~Moore, D.~Ramage, S.~Hampson, and B.~A. y Arcas,
\textit{Communication-efficient learning of deep networks from
decentralized data},
AISTATS, 2017.

\bibitem{Brinkmann19}
F.~Brinkmann, M.~Dinakaran, R.~Pelzer, P.~Grosche, D.~Voss, and S.~Weinzierl,
\textit{A cross-evaluated database of measured and simulated HRTFs
including 3D head meshes, anthropometric features, and headphone
impulse responses (HUTUBS)},
J.\ Audio Eng.\ Soc., 67(9):705--718, 2019.

\bibitem{Listen}
UMR 9912 - STMS - IRCAM/CNRS/UPMC,
\textit{LISTEN HRTF Database},
Olivier Warusfel.

\end{thebibliography}

\end{document}
